%% SCRAP: archive/quality/validation/phase-2-executive-summary
%% SOURCE: docs/working/archive/quality/validation/phase-2-executive-summary.md
%% STATUS: HISTORICAL
%% FITS: none
%% EDITORIAL: lifted — prose rewritten to press voice

\section{Phase 2 Experiment Executive Summary (November 2025)}

Dated 20 November 2025, this document captured the design for a coupled
multivariate dynamics experiment intended to supersede the Phase 1
factorial DoE. The experiment was designed but the Phase 1 baseline used
a fixed heartbeat; Phase 2 was to capture adaptive heartbeat behavior.
The work informed subsequent L8 attractor and window-scaling campaigns.

\subsection{The Paradigm Shift}

Phase 1 measured isolated factor effects under a fixed 1\,ms heartbeat
interval. Phase 2 recognized that the system is a coupled multivariate
dynamical system: the heartbeat rate (\texttt{tick\_ns}) is itself
modulated by the inference engine in response to workload, coupling all
metrics through a single feedback variable.

\subsection{Per-Tick Observables}

Seven metrics were to be captured per heartbeat tick via a circular
buffer (100\,K slots, approximately 400\,KB memory):

\begin{itemize}
  \item \texttt{tick\_interval\_ns} — elapsed time since previous tick
  \item \texttt{cache\_hits\_delta} — cache lookup successes
  \item \texttt{bucket\_hits\_delta} — secondary lookup successes
  \item \texttt{hot\_word\_count} — words above heat threshold
  \item \texttt{avg\_word\_heat} — mean execution heat
  \item \texttt{window\_width} — current rolling window size
  \item \texttt{word\_executions\_delta} — total word executions
\end{itemize}

\subsection{Experimental Design}

The design called for five elite configurations (identified from the
$2^7$ DoE), 150--200 runs each, across 100\,000+ ticks per run, for
750 total runs and an estimated 6--8 hour wall-clock runtime.

\begin{center}
\begin{tabular}{lp{6cm}}
\toprule
Dimension & Weight \\
\midrule
Stability   & 30\% — convergence and reproducibility \\
Convergence & 25\% — speed and reliability \\
Coupling    & 20\% — heartbeat response to load \\
Resilience  & 15\% — noise resistance \\
Jitter      & 10\% — timing cleanliness \\
\bottomrule
\end{tabular}
\end{center}

The configuration with the highest weighted composite score was to become
the baseline for MamaForth and PapaForth.

\subsection{Implementation Pipeline}

Five phases were planned: heartbeat instrumentation
(\texttt{HeartbeatTickSnapshot} struct and circular buffer injection),
per-run stability fingerprinting, per-configuration fingerprint aggregation,
golden configuration selection, and visualization with R/ggplot2.

%% TODO(bob): confirm whether the golden configuration was ever selected
%% and what label it carries in the current codebase. Cross-reference
%% against CLAUDE.md HeartbeatTickSnapshot note and
%% docs/03-architecture/heartbeat-system/instrumentation-plan.md.
